% Дополнение к главе 7. Основной файл должен определять \term, theorem и proof.
\subsection*{Корень произвольной натуральной степени}
\begin{definition}
\term{Корень нечётной степени} (odd-index root) \(n\geq 3\) из действительного числа \(a\) — единственное действительное число \(x\), для которого \(x^n=a\). Обозначение: \(x=\sqrt[n]{a}\).
\end{definition}
\begin{definition}
\term{Арифметический корень чётной степени} (principal even-index root) \(n\geq2\) из числа \(a\geq0\) — единственное \emph{неотрицательное} число \(x\), для которого \(x^n=a\). Обозначение: \(x=\sqrt[n]{a}\). При \(a<0\) такой действительный корень не определён.
\end{definition}
В общей записи \(\sqrt[n]{a}\) термин \term{корень \(n\)-й степени} (\(n\)th root) относится к одному из этих двух случаев. Здесь \(n\) — натуральное число, \(n\geq2\). Квадратный корень \(\sqrt a\) — это \(\sqrt[2]a\).

Не путайте значение корня с решениями уравнения: \(\sqrt[4]{16}=2\), а уравнение \(x^4=16\) имеет два действительных решения \(x=2\) и \(x=-2\). При нечётном \(n\) знак \(\sqrt[n]a\) совпадает со знаком \(a\).

\begin{theorem}[Корень из степени]
Для любого действительного \(x\) и натурального \(n\geq2\)
\[
\sqrt[n]{x^n}=
\begin{cases}
 x,&n\text{ нечётно},\\
 |x|,&n\text{ чётно}.
\end{cases}
\]
\end{theorem}
\begin{proof}
При нечётном \(n\) число \(x\) само является действительным корнем уравнения \(y^n=x^n\). При чётном \(n\) число \(|x|\) неотрицательно и \(|x|^n=x^n\), поэтому оно и является арифметическим корнем. Единственность в определении даёт оба равенства.
\end{proof}

\begin{theorem}[Произведение и частное корней]
Если \(n\geq2\) нечётно, то для любых действительных \(a,b\)
\[
\sqrt[n]{ab}=\sqrt[n]a\,\sqrt[n]b.
\]
Если \(n\) чётно, это равенство верно при \(a\geq0\) и \(b\geq0\). Для частного дополнительно требуется \(b\neq0\):
\[
\sqrt[n]{\frac ab}=\frac{\sqrt[n]a}{\sqrt[n]b}.
\]
Таким образом, при чётном \(n\) условия для частного: \(a\geq0\), \(b>0\); при нечётном: \(a\in\mathbb R\), \(b\in\mathbb R\setminus\{0\}\).
\end{theorem}
\begin{proof}
Пусть \(u=\sqrt[n]a\), \(v=\sqrt[n]b\). Тогда \((uv)^n=u^nv^n=ab\). При чётном \(n\) дополнительно \(uv\geq0\), поэтому \(uv\) — именно арифметический корень. Для частного \(v\neq0\) и \((\frac uv)^n=\frac ab\); при чётном \(n\) имеем \(\frac uv\geq0\). Снова применяем единственность корня.
\end{proof}

\begin{theorem}[Корень из корня]
Для натуральных \(m,n\geq2\) и \(a\geq0\)
\[
\sqrt[m]{\sqrt[n]a}=\sqrt[mn]a.
\]
Если \(m\) и \(n\) оба нечётны, равенство справедливо также при \(a<0\).
\end{theorem}
\begin{proof}
Обозначим левую часть через \(u\). Тогда \(u^m=\sqrt[n]a\), откуда \(u^{mn}=a\). При \(a\geq0\) число \(u\) неотрицательно; при отрицательном \(a\) и нечётных \(m,n\) степень \(mn\) нечётна. В каждом из указанных случаев \(u\) является корнем, определённым в правой части.
\end{proof}

\begin{example}
\(\sqrt[4]{80}=\sqrt[4]{16\cdot5}=2\sqrt[4]5\),
\(\sqrt[3]{-81}=-3\sqrt[3]3\),
\(\sqrt[6]{(-2)^6}=2\),
\(\sqrt[5]{(-2)^5}=-2\).
\end{example}

\subsection*{Рациональный показатель степени}
\begin{definition}
Пусть \(a>0\), \(m\in\mathbb Z\), \(n\in\mathbb N\). \term{Степень с рациональным показателем} (rational exponent) определяется равенством
\[
a^{\frac mn}=\bigl(\sqrt[n]a\bigr)^m.
\]
При \(n=1\) здесь имеется в виду обычная целая степень, а не новый корень. Отрицательный показатель означает обратную величину: \(a^{-\frac mn}=\frac{1}{a^{\frac mn}}\).
\end{definition}

\begin{theorem}[Независимость от записи дроби]
Для \(a>0\) значение \(a^{\frac mn}\) зависит только от рационального числа \(\frac mn\), а не от выбранных числителя и знаменателя.
\end{theorem}
\begin{proof}
Если \(\frac mn=\frac pq\), где \(n,q>0\), то \(mq=pn\). Оба числа \((\sqrt[n]a)^m\) и \((\sqrt[q]a)^p\) положительны. После возведения в натуральную степень \(nq\) они дают соответственно \(a^{mq}\) и \(a^{pn}\), то есть одинаковые числа. Два положительных числа с одинаковыми натуральными степенями равны. Последний факт следует из строгого возрастания натуральной степени на положительной полуоси: при \(0<u<v\) имеем \(u^N<v^N\).
\end{proof}

\begin{theorem}[Свойства рациональных степеней]
Пусть \(a,b>0\), \(r,s\in\mathbb Q\). Тогда
\[
\begin{aligned}
a^r a^s&=a^{r+s},&\qquad \frac{a^r}{a^s}&=a^{r-s},\\
(a^r)^s&=a^{rs},&(ab)^r&=a^rb^r,\\
\left(\frac ab\right)^r&=\frac{a^r}{b^r}.
\end{aligned}
\]
\end{theorem}
\begin{proof}
Приведём \(r=\frac mn\), \(s=\frac pn\) к общему положительному знаменателю и положим \(u=\sqrt[n]a>0\). Тогда \(a^ra^s=u^mu^p=u^{m+p}\); для частного рассуждение такое же. Для правила степени степени пусть \(r=\frac mn\), \(s=\frac pq\). Числа \((a^r)^s\) и \(a^{rs}\) положительны; их степени \(nq\) обе равны \(a^{mp}\), поэтому эти числа равны. Наконец, из свойства корня произведения получаем \((ab)^{\frac mn}=(\sqrt[n]a\sqrt[n]b)^m=a^{\frac mn}b^{\frac mn}\); для частного аналогично. Целые отрицательные показатели допускаются, поскольку все основания и корни ненулевые.
\end{proof}

\subsubsection*{Нулевое и отрицательное основания}
Для \(r>0\) принимают \(0^r=0\). Записи \(0^0\) и \(0^r\) при \(r<0\) в этом сборнике не определены.

Для \(a<0\) используем следующую договорённость: сначала записываем рациональное число \(r=\frac mn\) как \term{несократимую дробь} (reduced fraction), где \(n>0\). Действительную степень \(a^{\frac mn}\) определяем только при нечётном \(n\):
\[
a^{\frac mn}=\bigl(\sqrt[n]a\bigr)^m.
\]
При чётном знаменателе несократимой дроби такая действительная степень не определена. Например,
\[
(-243)^{\frac25}=9,\qquad
(-27)^{-\frac13}=-\frac13,\qquad
(-4)^{\frac12}\text{ не определена в }\mathbb R.
\]
Общими правилами степеней для отрицательных оснований пользоваться без проверки нельзя. Например,
\[
\bigl((-1)^2\bigr)^{\frac12}=1,
\qquad
(-1)^{2\cdot(\frac12)}=-1.
\]
Также \((-8)^{\frac26}=(-8)^{\frac13}=-2\), тогда как \(\sqrt[6]{(-8)^2}=2\). Условие \(a>0\) в теореме о свойствах рациональных степеней существенно.
